
Or il
résulte de 2.2 que l'isomorphisme:
@@M0@@
défini par les isomorphismes (2.5.2) est compatible aux structures d'algèbres des deux membres; il en est donc de même pour l'isomorphisme:
@@M1@@
défini par les isomorphismes (2.5.3). Pour des raisons d'unicité, ce dernier est donc la restriction de l'homomorphisme (2.5.1) qui avait été défini globalement. L'homomorphisme (2.5.1) est donc un isomorphisme.

ii) On considère la suite exacte:
@@M2@@
et la suite exacte des @@M3@@ qui en résulte:
@@M4@@
Il en résulte que @@M5@@ est le noyau de l'homomorphisme :
@@M6@@
c'est-à-dire @@M7@@.

Pour montrer la dernière relation, nous allons montrer qu'il existe une suite exacte:
@@M8@@
Les homomorphismes @@M9@@ et @@M10@@ ont été définis en (1.6.1); on prend ici leur différence. Par ailleurs, @@M11@@ s'envoie de façon naturelle dans @@M12@@, donc dans @@M13@@, et de même dans @@M14@@, et ces homomorphismes s'annulent sur @@M15@@, car @@M16@@ (où @@M17@@ désigne l'idéal de @@M18@@) ce qui définit la flèche de gauche.

